\documentclass{article}

\usepackage[main, preprint]{neurips_2025}

\usepackage[utf8]{inputenc}
\usepackage[T1]{fontenc}
\usepackage{hyperref}
\usepackage{url}
\usepackage{graphicx}
\usepackage{booktabs}
\usepackage{amsfonts}
\usepackage{amsmath}
\usepackage{amssymb}
\usepackage{nicefrac}
\usepackage{microtype}
\usepackage{xcolor}
\usepackage{algorithm}
\usepackage{algorithmic}
\usepackage{subcaption}

\title{Stylized Rendering as a Function of Expectation:\\
  Implementation and Extension in La~Jolla}

\author{%
  Yuezhexuan Zhu \\
  Department of Computer Science and Engineering\\
  University of California, San Diego\\
  \texttt{jaz019@ucsd.edu} \\
}

\begin{document}

\maketitle

%% ═══════════════════════════════════════════════════════════════
%%  ABSTRACT
%% ═══════════════════════════════════════════════════════════════
\begin{abstract}
We present an implementation and extension of West's~\cite{west2024stylized} framework for stylized rendering as a function of expectation within the La~Jolla physically-based renderer. Starting from a deterministic cel-shading pipeline (\textbf{Pipeline~1}), we progressively build toward a Monte~Carlo--based stylization system that applies non-linear style functions $g_\theta$ to per-pixel radiance estimates (\textbf{Pipeline~2}). We then add object-selective stylization with temporal crossfade (\textbf{Pipeline~3}) and culminate in a \textbf{multi-style showcase} (\textbf{Pipeline~4}) featuring per-object dispatch of three distinct style functions---cel-shading, tie-dye cosine mapping, and artistic colour palette (ACP)---demonstrated in a custom open-air \emph{Sculpture Garden} scene with animated lighting. Our results demonstrate that stylization applied to the \emph{expectation} of Monte~Carlo transport yields aesthetically coherent effects with physically grounded illumination, including global illumination, soft shadows, and colour bleeding.
\end{abstract}


%% ═══════════════════════════════════════════════════════════════
%%  1. INTRODUCTION
%% ═══════════════════════════════════════════════════════════════
\section{Introduction}
\label{sec:intro}

Non-photorealistic rendering (NPR) has a long history in computer graphics, from Gooch shading~\cite{gooch1998non} and toon/cel-shading to more recent painterly and hatching styles. However, most traditional NPR techniques operate in a simplified shading model (e.g.\ a single directional light, no indirect illumination) or are applied as image-space post-processes that are blind to 3D light transport.

West~\cite{west2024stylized} proposes an elegant alternative: apply a \emph{non-linear style function} $g_\theta$ to the Monte~Carlo \emph{expectation} of physically based radiance at each pixel. Because the style function operates on the converged estimate $\langle I \rangle$ rather than individual samples, the resulting image inherits full global illumination (GI) effects—soft shadows, colour bleeding, caustics—while exhibiting the desired artistic look.

In this project, we implement and extend this framework within the La~Jolla renderer~\cite{li2018differentiable} across four progressively more capable pipelines:

\begin{enumerate}
    \item \textbf{Pipeline~1} (Skypop NPR): Deterministic cel-shading with Sobel-based outlines. Single directional light, zero bounces.
    \item \textbf{Pipeline~2} (West Static): $k$-sample Monte~Carlo estimate $\langle I \rangle$ with a cel-shading $g_\theta$. Full GI.
    \item \textbf{Pipeline~3} (West Animated): Object-ID--selective stylization with temporal crossfade $t \in [0,1]$, animated camera orbit.
    \item \textbf{Pipeline~4} (Multi-Style Showcase): Per-object dispatch of three style functions (cel, tie-dye, ACP) in a custom \emph{Sculpture Garden} scene with a sweeping area light.
\end{enumerate}

Our primary contributions are: (a)~a clean implementation of West's stylization framework in La~Jolla, (b)~a per-object style dispatch mechanism enabling heterogeneous styles in a single render, and (c)~a custom open-air scene with animated lighting that demonstrates the interplay between physically based illumination and stylistic colour transformations.


%% ═══════════════════════════════════════════════════════════════
%%  2. RELATED WORK
%% ═══════════════════════════════════════════════════════════════
\section{Related Work}
\label{sec:related}

\paragraph{Traditional NPR.}
Gooch et al.~\cite{gooch1998non} introduced a warm-to-cool shading model that maps $\mathbf{N} \cdot \mathbf{L}$ to a colour ramp, providing shape cues without photorealism. Lake et al.~\cite{lake2000stylized} surveyed cartoon-style rendering, including silhouette extraction and quantised shading. These methods typically use a single light source and ignore indirect illumination.

\paragraph{Stylized path tracing.}
West~\cite{west2024stylized} formalises the idea of applying a style function $g_\theta$ to the \emph{expected} exitant radiance at each visible surface point. The key insight is that the non-linear function should be applied \emph{after} averaging multiple Monte~Carlo samples, preserving the convergence properties of the estimator while introducing stylistic non-linearity. This differs from naively stylizing each sample independently, which produces noisy results that do not converge to the desired style.

\paragraph{Artistic colour palettes.}
Colour mapping and palette-based rendering have been explored in the context of image stylization~\cite{reinhard2001color}. West~\cite{west2024stylized} demonstrates two such mappings: a per-channel cosine ``tie-dye'' effect and a luminance-based Artistic Colour Palette (ACP) ramp, both of which we implement in Pipeline~4.

% TODO: Add more related work as needed (e.g., halftone rendering, painterly rendering)


%% ═══════════════════════════════════════════════════════════════
%%  3. METHOD
%% ═══════════════════════════════════════════════════════════════
\section{Method}
\label{sec:method}

\subsection{Background: Monte Carlo Rendering}
\label{sec:mc}

In physically based rendering, the colour of each pixel is the integral of the measurement function over the pixel's area and the light transport paths. The rendering equation~\cite{kajiya1986rendering} defines exitant radiance $L_o$ at a surface point $\mathbf{x}$ in direction $\omega_o$ as:
\begin{equation}
    L_o(\mathbf{x}, \omega_o) = L_e(\mathbf{x}, \omega_o) + \int_{\mathcal{H}^2} f_s(\mathbf{x}, \omega_i, \omega_o)\, L_i(\mathbf{x}, \omega_i)\, |\cos\theta_i|\, \mathrm{d}\omega_i
    \label{eq:rendering}
\end{equation}
where $L_e$ is emission, $f_s$ is the BSDF, and the integral is over the hemisphere $\mathcal{H}^2$. Monte~Carlo path tracing estimates this integral by averaging $N$ random samples:
\begin{equation}
    \langle I \rangle = \frac{1}{N} \sum_{j=1}^{N} \frac{f_s \cdot L_i \cdot |\cos\theta_i|}{p(\omega_j)}
    \label{eq:mc_estimate}
\end{equation}


\subsection{Stylized Rendering Framework}
\label{sec:framework}

Following West~\cite{west2024stylized}, we define stylized rendering as the application of a non-linear style function $g_\theta$ to the expected radiance at each visible surface point. For a camera ray hitting surface point $\mathbf{x}$ with material albedo $\rho$, the stylized pixel colour is:
\begin{equation}
    C_{\text{styled}} = g_\theta\!\bigl(\rho,\; \langle I \rangle_k\bigr)
    \label{eq:stylized}
\end{equation}
where $\langle I \rangle_k$ is the $k$-sample Monte~Carlo estimate of the physical radiance at $\mathbf{x}$. The key property is that $g_\theta$ is applied to the \emph{average} of $k$ inner samples, not to individual samples.

\paragraph{Inner trace.} Each of the $k$ inner samples is a full path trace from the first camera-hit vertex, using next-event estimation (NEE) with MIS weighting and Russian roulette for termination:
\begin{equation}
    \langle I \rangle_k = \frac{1}{k} \sum_{s=1}^{k} I_s, \qquad I_s = \sum_{b=0}^{B} T_b \cdot \bigl(w_1 \cdot C_{\text{NEE}}^{(b)} + w_2 \cdot C_{\text{BSDF}}^{(b)}\bigr)
    \label{eq:inner_trace}
\end{equation}
where $T_b$ is the path throughput at bounce $b$, $w_1, w_2$ are MIS weights, and $B$ is the path length (controlled by Russian roulette).


\subsection{Style Functions $g_\theta$}
\label{sec:style_functions}

We implement three style functions, each corresponding to a different visual effect from West~\cite{west2024stylized}.

\subsubsection{Cel-Shading (Pipeline 1--4)}
\label{sec:cel}

A hard-threshold quantisation of luminance into lit and shadow bands:
\begin{equation}
    g_{\text{cel}}(\rho, \langle I \rangle) = 
    \begin{cases}
        \rho \cdot \mathbf{c}_{\text{light}} + \rho \cdot \mathbf{c}_{\text{ambient}} & \text{if } \text{lum}(\langle I \rangle) \geq \tau \\
        \rho \cdot \mathbf{c}_{\text{shadow}} + \rho \cdot \mathbf{c}_{\text{ambient}} & \text{otherwise}
    \end{cases}
    \label{eq:cel}
\end{equation}
where $\tau$ is the cel threshold, $\mathbf{c}_{\text{light}}$ and $\mathbf{c}_{\text{shadow}}$ are the light and shadow tint colours, and $\text{lum}(\cdot)$ computes the luminance of a colour vector.

\paragraph{Mollification effect.} At low $k$, variance in $\langle I \rangle$ causes pixels near $\tau$ to stochastically flip between lit and shadow, producing a soft, anti-aliased boundary. At high $k$, the estimate converges and the boundary becomes sharp. This is the super-linear convergence property described in West~\cite{west2024stylized} \S3.2.

\subsubsection{Tie-Dye Cosine Effect (Pipeline 4)}
\label{sec:tiedye}

A psychedelic colour mapping using per-channel cosine waves (West~\cite{west2024stylized}, Fig.~11):
\begin{equation}
    g_{\text{tiedye}}(\langle I \rangle)_c = \bigl|\cos\!\bigl(\nu_c \cdot \langle I \rangle_c + \phi_c\bigr)\bigr|, \qquad c \in \{r, g, b\}
    \label{eq:tiedye}
\end{equation}
where $\nu_c$ and $\phi_c$ are the per-channel frequency and phase parameters. Because the cosine operates on the physical radiance, the resulting colour pattern is modulated by global illumination—shadows, reflections, and colour bleeding all produce different hues.

\subsubsection{Artistic Colour Palette / ACP (Pipeline 4)}
\label{sec:acp}

A luminance-to-colour-ramp mapping (West~\cite{west2024stylized}, Fig.~7):
\begin{equation}
    g_{\text{acp}}(\langle I \rangle) = \text{lerp}\!\bigl(\mathbf{c}_{\text{dark}},\; \mathbf{c}_{\text{bright}},\; \text{clamp}\bigl(\text{lum}(\langle I \rangle), 0, 1\bigr)\bigr)
    \label{eq:acp}
\end{equation}
This maps all pixels to a two-stop gradient, creating a limited-palette illustration effect. Our implementation uses $\mathbf{c}_{\text{dark}} = (0.06, 0.02, 0.18)$ (deep indigo) and $\mathbf{c}_{\text{bright}} = (1.0, 0.50, 0.15)$ (sunset orange).


\subsection{Per-Object Style Dispatch}
\label{sec:dispatch}

A key extension beyond West's original framework is our per-object style dispatch mechanism. Rather than applying a single $g_\theta$ to all objects, we assign each shape a style via three separate ID lists parsed from the scene XML:

\begin{equation}
    \text{style}(\texttt{shape\_id}) = 
    \begin{cases}
        \text{CEL}    & \text{if } \texttt{shape\_id} \in \mathcal{S}_{\text{cel}} \\
        \text{TIEDYE} & \text{if } \texttt{shape\_id} \in \mathcal{S}_{\text{tiedye}} \\
        \text{ACP}    & \text{if } \texttt{shape\_id} \in \mathcal{S}_{\text{acp}} \\
        \text{PBR}    & \text{otherwise}
    \end{cases}
    \label{eq:dispatch}
\end{equation}

Non-target objects skip the expensive $k$-sample inner loop entirely, using a single-sample PBR estimate instead, yielding approximately $k\!\times$ speedup for those pixels.


\subsection{Temporal Crossfade}
\label{sec:crossfade}

For animated transitions, we introduce a parameter $t \in [0, 1]$ that blends between the physical PBR estimate and the stylized output:
\begin{equation}
    C_{\text{final}} = (1 - t) \cdot \langle I \rangle_k + t \cdot g_\theta(\rho, \langle I \rangle_k)
    \label{eq:crossfade}
\end{equation}
At $t = 0$ the image is fully photorealistic; at $t = 1$ it is fully stylized. Intermediate values produce a smooth blend.


%% ═══════════════════════════════════════════════════════════════
%%  4. IMPLEMENTATION
%% ═══════════════════════════════════════════════════════════════
\section{Implementation}
\label{sec:implementation}

\subsection{Renderer Architecture}

All pipelines are implemented within the La~Jolla physically-based renderer~\cite{li2018differentiable}, a C++ path tracer with Embree-based ray-scene intersection. The key modified/added source files are:

\begin{itemize}
    \item \texttt{stylized\_pt\_integrator.h}: Core integrator with inner trace, style functions, per-object dispatch, and temporal crossfade.
    \item \texttt{npr\_integrator.h}: Deterministic cel-shading integrator (Pipeline~1) with AOV output.
    \item \texttt{scene.h}: Extended \texttt{RenderOptions} struct with per-style ID lists and style parameters.
    \item \texttt{parse\_scene.cpp}: XML parsing for \texttt{cel\_ids}, \texttt{tiedye\_ids}, \texttt{acp\_ids}, and associated parameters.
\end{itemize}

\subsection{Pipeline~1: Deterministic Cel-Shading}

Pipeline~1 uses a single directional light and computes $\mathbf{N} \cdot \mathbf{L}$ with a hard shadow ray. A Sobel edge detection post-process on depth, normal, and object-ID AOVs produces outlines (see Section~\ref{sec:postproc}). This serves as our NPR baseline with zero Monte~Carlo noise.

\subsection{Pipeline~2--3: Monte Carlo Stylization}

Pipelines~2 and~3 use the $k$-sample inner trace (Eq.~\ref{eq:inner_trace}) with full path tracing: NEE, BSDF sampling, MIS, and Russian roulette. Pipeline~3 adds object-ID filtering (Eq.~\ref{eq:dispatch}) and the temporal crossfade (Eq.~\ref{eq:crossfade}).

\subsection{Pipeline~4: Multi-Style Showcase}

Pipeline~4 extends the dispatch to three simultaneous style functions and introduces a custom \emph{Sculpture Garden} scene (Section~\ref{sec:scene}) with animated lighting.

\subsection{Sculpture Garden Scene}
\label{sec:scene}

To demonstrate the multi-style system beyond the Cornell Box, we designed a custom open-air scene with procedurally generated OBJ meshes:

\begin{itemize}
    \item A 60$\times$60 ground plane (warm sandy stone)
    \item Three pedestal variants (tall, short, wide) for displaying objects
    \item A 32-sided cylindrical column
    \item Four spheres at various heights
    \item A large overhead area light (positioned above the camera frustum to remain invisible)
\end{itemize}

Each shape is assigned a style: golden pedestals and sky-blue/pearl spheres receive \textbf{cel-shading}, the emerald column and lavender orb receive \textbf{tie-dye}, and terracotta/coral/slate objects receive \textbf{ACP}. The camera is fixed while the area light sweeps in a semicircular arc from left-grazing to overhead to right-grazing, demonstrating how each style responds differently to changing illumination geometry.

\subsection{Animation Pipeline}

The animation script (\texttt{animate\_gallery.py}) generates 60 frames by patching the area light's position in the scene XML (note: La~Jolla resolves OBJ mesh paths relative to the scene XML directory, so temporary frame XMLs must be written in the scene directory):
\begin{equation}
    \mathbf{p}_{\text{light}}(i) = \begin{pmatrix} -R\cos(\pi i / (N-1)) \\ y_{\min} + (y_{\max} - y_{\min})\sin(\pi i / (N-1)) \\ z_0 \end{pmatrix}
    \label{eq:light_sweep}
\end{equation}
where $R$ is the sweep radius, $y_{\min}$ and $y_{\max}$ define the height range, and $N = 60$ frames. Frames are rendered in parallel, converted from EXR to PNG via \texttt{ffmpeg}, and assembled into an animated GIF using palette-based dithering for quality.


%% ═══════════════════════════════════════════════════════════════
%%  5. NPR POST-PROCESSING
%% ═══════════════════════════════════════════════════════════════
\section{NPR Post-Processing Baseline}
\label{sec:postproc}

Before implementing Monte~Carlo stylisation (Pipelines~2--4), we developed a suite of 2D NPR post-processing techniques as a baseline. These operate on finished images and serve as a point of comparison for the 3D-integrated approach.

\subsection{Ordered Dithering (Obra Dinn Style)}

Inspired by \emph{Return of the Obra Dinn}~\cite{obradinn2018}, we convert any image to a two-colour palette via Bayer ordered dithering~\cite{floyd1976adaptive}. The algorithm: (1)~tone-map to LDR, (2)~extract luminance, (3)~threshold against a tiled $n \times n$ Bayer matrix ($n \in \{2,4\}$, yielding 4$\times$4 to 16$\times$16 patterns), and (4)~remap the binary mask to a chosen palette (sepia, green-phosphor, etc.). Optional Sobel edge enhancement before dithering preserves silhouettes.

\subsection{Toon / Cel Quantisation}

Screen-space cel shading quantises the luminance channel into $B$ discrete bands:
\begin{equation}
    L_{\text{cel}}(x) = \left\lfloor L(x) \cdot B \right\rfloor / B
    \label{eq:cel_quant}
\end{equation}
where $L(x)$ is the pixel luminance. Steps: (1)~convert to HSL, (2)~snap lightness to $B$ levels, (3)~recompose HSL $\to$ RGB (preserving hue and saturation), (4)~detect edges via 4-connected luminance differences and paint them black. This is the image-space analogue of the 3D cel-shading in Pipeline~1.

\subsection{Painterly Rendering}

Three painting styles are built on the Kuwahara anisotropic smoothing filter~\cite{kuwahara1976processing}, which replaces each pixel with the mean of the lowest-variance quadrant in its neighbourhood, producing characteristic oil-paint blockiness. The \emph{Litwinowicz} variant~\cite{litwinowicz1997processing} additionally overlays oriented brush strokes along image-gradient directions. Colour palette reduction (uniform or $k$-means) limits the output to a target palette size.

\subsection{Sobel Outline Post-Processing}
\label{sec:sobel}

For Pipeline~1, outline edges are extracted from three AOV buffers (depth, surface normal, object~ID) using a G-buffer--based approach~\cite{saito1990comprehensible}. Rather than a full Sobel convolution, we use direct 4-connected neighbour comparison:
\begin{equation}
    \text{edge}(x) = \bigl(|d_x - d_{\text{nbr}}| > \tau_d\bigr) \;\lor\; \bigl(\|\mathbf{n}_x - \mathbf{n}_{\text{nbr}}\| > \tau_n\bigr) \;\lor\; \bigl(\text{id}_x \neq \text{id}_{\text{nbr}}\bigr)
    \label{eq:sobel}
\end{equation}
Silhouette pixels (geometry vs.\ background boundary) are always marked. A camera-facing normal flip ($\mathbf{N} \to -\mathbf{N}$ when $\mathbf{N} \cdot (-\mathbf{d}) < 0$) resolved a winding-order artifact on the Cornell Box ceiling, reducing outlined pixels by 35\%.


%% ═══════════════════════════════════════════════════════════════
%%  6. RESULTS
%% ═══════════════════════════════════════════════════════════════
\section{Results}
\label{sec:results}

% TODO: Add rendered images here. Suggested figures:
% - Figure 1: Pipeline 1 vs Pipeline 2 comparison (deterministic vs MC cel-shading)
% - Figure 2: Pipeline 3 temporal crossfade sequence (t=0, 0.33, 0.67, 1.0)
% - Figure 3: Pipeline 4 Sculpture Garden — cel/tiedye/acp side by side
% - Figure 4: Light sweep animation frames (left-grazing, overhead, right-grazing)
% - Table 1: Rendering time comparison across pipelines and configurations

\subsection{Visual Comparison}

\begin{figure}[h]
    \centering
    % \includegraphics[width=\linewidth]{figures/pipeline_comparison.png}
    \fbox{\rule[-.5cm]{0cm}{4cm} \rule[-.5cm]{14cm}{0cm}}
    \caption{Comparison across pipelines. \textbf{Left}: Pipeline~1 (deterministic cel + Sobel outlines). \textbf{Centre}: Pipeline~2 (MC cel, $k = 256$). \textbf{Right}: Pipeline~4 (multi-style: cel sphere, tie-dye column, ACP pedestal). Pipeline~2--4 exhibit physically correct soft shadows and colour bleeding absent in Pipeline~1.}
    \label{fig:comparison}
\end{figure}

% TODO: Replace placeholder with actual rendered images

\subsection{Style Function Effects}

% TODO: Add close-up renders showing each style function's effect on the same object
% under different lighting conditions

\begin{figure}[h]
    \centering
    % \includegraphics[width=\linewidth]{figures/style_functions.png}
    \fbox{\rule[-.5cm]{0cm}{3cm} \rule[-.5cm]{14cm}{0cm}}
    \caption{Style function comparison on a sphere under identical lighting. \textbf{(a)}~PBR baseline. \textbf{(b)}~Cel-shading: hard lit/shadow boundary. \textbf{(c)}~Tie-dye: per-channel cosine waves produce psychedelic colour banding. \textbf{(d)}~ACP: luminance-to-gradient mapping (indigo $\to$ sunset orange).}
    \label{fig:styles}
\end{figure}

\subsection{Light Sweep Animation}

% TODO: Add key frames from the animated GIF showing the light sweeping
% from left-grazing → overhead → right-grazing

\begin{figure}[h]
    \centering
    % \includegraphics[width=\linewidth]{figures/light_sweep.png}
    \fbox{\rule[-.5cm]{0cm}{3cm} \rule[-.5cm]{14cm}{0cm}}
    \caption{Selected frames from the light-sweep animation (60 frames, 15~fps). \textbf{Left}: Grazing illumination from the left ($i = 0$). \textbf{Centre}: Overhead lighting ($i = 30$). \textbf{Right}: Grazing from the right ($i = 60$). Each style reacts distinctly to the lighting change: cel boundaries shift, tie-dye bands modulate, and ACP gradients rotate along the colour ramp.}
    \label{fig:light_sweep}
\end{figure}

\subsection{Performance}

\begin{table}[h]
    \caption{Rendering time per frame at 640$\times$480 resolution on an Apple M3 14-core CPU.}
    \label{tab:performance}
    \centering
    \begin{tabular}{lcccc}
        \toprule
        Pipeline & SPP & Inner $k$ & Style & Time (s) \\
        \midrule
        1 (Skypop NPR) & 1 & --- & Cel (deterministic) & $\sim$0.05 \\
        2 (West Static) & 4 & 128 & Cel & $\sim$0.45 \\
        3 (West Animated) & 4 & 128 & Cel (selective) & $\sim$0.35 \\
        4 (Multi-Style) & 16 & 256 & Cel/Tiedye/ACP & $\sim$2.8 \\
        \bottomrule
    \end{tabular}
\end{table}

Pipeline~1 is effectively instantaneous due to its deterministic single-ray evaluation. Pipelines~2--3 scale linearly with $k$. Pipeline~4 uses higher $k$ (256) and SPP (16) for quality renders; non-target objects remain cheap (1-sample PBR fallback).

% TODO: Add convergence analysis — how cel boundaries sharpen with increasing k


%% ═══════════════════════════════════════════════════════════════
%%  6. DISCUSSION & FUTURE WORK
%% ═══════════════════════════════════════════════════════════════
\section{Discussion and Future Work}
\label{sec:future}

\paragraph{Convergence properties.}
The super-linear convergence of the cel threshold is readily apparent: at $k = 16$ the shadow boundary is soft and noisy, while at $k = 256$ it is crisp and deterministic. This mollification effect, described in West~\cite{west2024stylized} \S3.2, is a direct consequence of applying $g_\theta$ to the averaged estimate rather than individual samples.

\paragraph{Limitations.}
Our implementation has several limitations:
\begin{itemize}
    \item The ACP colour ramp is limited to a two-stop linear gradient; a multi-stop or lookup-table-based mapping would provide richer palettes.
    \item The tie-dye frequency/phase parameters are global; per-object parameterisation would allow more diverse effects.
    \item Cross-hatching and other geometry-dependent styles (West~\cite{west2024stylized}, \S4) were attempted but abandoned due to shadow-ray self-intersection issues with procedural 3D hatching planes.
\end{itemize}

\paragraph{Future directions.}
\begin{enumerate}
    \item \textbf{Multi-stop ACP}: Extend the ACP ramp to $n$-stop colour gradients or artist-designed lookup textures.
    \item \textbf{Procedural hatching}: Revisit cross-hatching using object-space parameterisation and tangent-frame alignment to avoid the planar projection artefacts encountered in our initial attempt.
    \item \textbf{Temporal coherence}: Apply the style function to temporally smoothed radiance estimates for flicker-free animation.
    \item \textbf{GPU acceleration}: Port the inner trace loop to CUDA/OptiX for interactive preview rates.
    \item \textbf{Perceptual evaluation}: Conduct user studies to assess the aesthetic quality and 3D shape perception of different style functions under varying illumination.
\end{enumerate}


%% ═══════════════════════════════════════════════════════════════
%%  7. CONCLUSION
%% ═══════════════════════════════════════════════════════════════
\section{Conclusion}
\label{sec:conclusion}

We have implemented and extended West's framework for stylized rendering as a function of expectation within the La~Jolla renderer. Our four-pipeline progression demonstrates the power of applying non-linear style functions to Monte~Carlo radiance estimates: the resulting images combine physically correct global illumination with artistically coherent stylization. The per-object dispatch mechanism and multi-style Sculpture Garden scene showcase the flexibility of this approach for complex, heterogeneous scenes. Code and scenes are available at \url{https://github.com/zhuyuezx/lajolla_public}.


%% ═══════════════════════════════════════════════════════════════
%%  REFERENCES
%% ═══════════════════════════════════════════════════════════════
\bibliographystyle{unsrt}
\bibliography{references}

%%%%%%%%%%%%%%%%%%%%%%%%%%%%%%%%%%%%%%%%%%%%%%%%%%%%%%%%%%%%

\appendix

\section{Integrator Pseudocode}
\label{app:pseudocode}

\begin{algorithm}[h]
\caption{Stylized Path Tracing with Per-Object Dispatch}
\label{alg:stylized_pt}
\begin{algorithmic}[1]
\REQUIRE Scene $\mathcal{S}$, ray $(o, d)$, params $\theta$
\STATE $\mathbf{x} \leftarrow \text{Intersect}(\mathcal{S}, o, d)$
\IF{miss}
    \RETURN $\mathbf{c}_{\text{bg}}$
\ENDIF
\STATE $\text{style} \leftarrow \text{GetStyleForShape}(\texttt{shape\_id}(\mathbf{x}))$ \hfill {\footnotesize // Eq.~\ref{eq:dispatch}}
\IF{$\text{style} = \text{PBR}$}
    \RETURN $\text{InnerTrace}(\mathbf{x}, 1)$ \hfill {\footnotesize // single sample, fast}
\ENDIF
\STATE $\langle I \rangle \leftarrow \frac{1}{k}\sum_{s=1}^{k} \text{InnerTrace}(\mathbf{x}, s)$ \hfill {\footnotesize // Eq.~\ref{eq:inner_trace}}
\STATE $C_{\text{styled}} \leftarrow g_\theta(\rho(\mathbf{x}),\, \langle I \rangle)$ \hfill {\footnotesize // Eq.~\ref{eq:cel}/\ref{eq:tiedye}/\ref{eq:acp}}
\RETURN $(1 - t) \cdot \langle I \rangle + t \cdot C_{\text{styled}}$ \hfill {\footnotesize // Eq.~\ref{eq:crossfade}}
\end{algorithmic}
\end{algorithm}


\section{Scene Configuration (XML)}
\label{app:xml}

% TODO: Add example XML snippets showing per-object style assignment


\end{document}